\section{What remains, ranked}\label{sec:outlook}

The decomposition closes some questions and opens others. We list what is
closed first, so that it is not redone, and then what is open, in the order we
would take it.

\subsection{Closed}

\begin{enumerate}
\item The comb is not a Loewner phenomenon, in either identification of the
variables, for single- or multi-tip chains, deterministic or L\'evy-driven
(Section~\ref{sec:nocomb}).
\item The rational correction carries no arithmetic and its residues are not
data (Proposition~\ref{prop:forced}); the inherited note's expectation that a
realization would fail ``at two numbers'' is withdrawn.
\item Structural positivity of the Markov part bounds nothing about
$\norm{V_{\omega,L}}$ (Section~\ref{sec:decomposition}); the cancellation
between the positive part and its EMA is the content.
\item The far-right-line Fourier realization of the transfer is superseded for
numerical purposes by the elementary assembly (Section~\ref{sec:assembly}).
\end{enumerate}

\subsection{Open}

\begin{enumerate}
\item \textbf{Finish the contraction side.} A standard-library registered
version of the assembly at $L=\log3$, the one recorded horizon where
$2\omega m_L\approx10^{-9}$ is resolvable in double precision, asserting
contraction, $\lambda_{\min}(D_N)/2\omega=m_L^{(N)}(1+\delta)$ with $|\delta|<1\%$
at $\omega=0.01$, and the mass against Proposition~\ref{prop:tail}. The
horizons $\log5,\log7$ at 40 digits, exploratory. The cumulative Cayley
coordinate of \cite[Proposition~7.7]{WilsonLines} from the same matrices. And
the phase-matched Gram control of \cite{WilsonReview} on the form side against
the transfer's margin: the transfer has no lattice analogue, so this control
measures what the transfer knows that a lattice does not.
\item \textbf{The Lax--Phillips dictionary at $\omega=\frac12$}, object by
object: which subspace the compression $P_L$ is in their energy form, what the
flux identity \eqref{eq:flux} is there, and what their semigroup estimate says
about the contraction margin at the endpoint. Then which of their objects the
$\omega$-deformation moves and which it cannot reach.
\item \textbf{The Beta law from a radial chain.} Whether a conformal-radius or
hitting law of radial Loewner evolution driven by a Bessel process of dimension
$b$ has the additivity fraction of Proposition~\ref{prop:bessel} as a squared
radius ratio. If yes, the archimedean factor has a Loewner realization and its
quarter is a Bessel dimension; if no, the archimedean factor is a
Bessel-process statement and not a Loewner one, and the investigation's name
records its origin rather than a direction.
\item \textbf{The comb in the Bost--Connes algebra.} $C_\omega$ is the image of
$\zeta(s-\omega)/\zeta(s+\omega)$ under $n^{-s}\mapsto\mu_n$. What the KMS
structure of that algebra says about the compression $P_LC_\omega P_L$, whose
norm grows like $e^{bL}$, and whether the subtraction $2\EMA_b$ has a meaning
there. A reading task with a small chance of a real statement.
\item \textbf{The criticality in filter terms.} The transfer is a lossless
filter with impulse response $\kh_\omega-2\EMA_b[\kh_\omega]$, $\kh_\omega\geq0$,
and RH is its stability at every $\omega>0$. A clean statement of why no
passivity theorem for that class can avoid the location of the poles --- what
property of $\kh_\omega$ beyond positivity would be needed --- would be the
honest end of this line.
\end{enumerate}

\subsection{Status of every statement}

\emph{Written proof, unconditional:} Propositions~\ref{prop:unimodular},
\ref{prop:factor}--\ref{prop:blaschke}, \ref{prop:bessel}, \ref{prop:farfield},
\ref{prop:secondorder}--\ref{prop:tail}, Lemmas~\ref{lem:multiplicative}
and~\ref{lem:rotation}, and the identities of Section~\ref{sec:endpoint}.
\emph{Quoted from \cite{WilsonLines}:} Proposition~\ref{prop:dichotomy} and the
flow identities. \emph{Labelled numerical computation (unregistered):}
Section~\ref{sec:results} and Appendix~\ref{app:numerics}. \emph{Reading:} the
Lax--Phillips remark of Section~\ref{sec:endpoint}; the Bost--Connes
identification beyond the algebra written out; the $1/N'$ law and the two-point
fit of Section~\ref{sec:results}. \emph{Not claimed:} a realization of the
transfer, a contraction at any horizon beyond $\log3$, a positivity statement,
convergence of the sine basis, or anything about the zeros.
